\makeatletter
\def\input@path{{./}{../}{../../}{../../../}{../../../../}{preamble/}{../preamble/}{../../preamble/}{../../../preamble/}}
\makeatother

\input{preamble_exams.tex}

\title{Decision-Making Under Uncertainty: Exam Solutions}
\author{}
\date{}

\begin{document}
	\maketitle
	\vspace{-2cm}
	\setcounter{tocdepth}{1}
	\setcounter{secnumdepth}{2}
	\phantomsection
	\label{toc}
	\tableofcontents

	\vspace{-5mm}
	\section*{Learning objective}\vspace{-5mm}
	Apply minimax regret and Expected Monetary Value (EMV) to select an alternative under uncertainty, and use sensitivity analysis to determine how a subjective probability changes the preferred alternative.

		\section*{Criteria assessment}\vspace{-5mm}
		Each problem is worth 2 points. Show all calculations, including the complete regret table for C5, every EMV calculation for C6, and the threshold equation and decision ranges for C7. Partial credit may be awarded when the correct procedure is clearly shown but arithmetic errors are made. No points will be awarded if no procedure is shown.
		A \textbf{minimum of 4 points} is required to pass each criterion. A 1-point penalty will be applied to each criterion if the student's name is not written on the assessment or if the assessment is submitted late.

		\begin{center}
			\begin{tabular}{|c|c|c|c|}
				\hline
				\textbf{Question} & \textbf{C5} & \textbf{C6} & \textbf{C7} \\ \hline
				1 & 2 & -- & -- \\ \hline
				2 & 2 & -- & -- \\ \hline
				3 & 2 & -- & -- \\ \hline
				4 & -- & 2 & -- \\ \hline
				5 & -- & 2 & -- \\ \hline
				6 & -- & 2 & -- \\ \hline
				7 & -- & -- & 2 \\ \hline
				8 & -- & -- & 2 \\ \hline
				9 & -- & -- & 2 \\ \hline
				\textbf{Total} & \textbf{6} & \textbf{6} & \textbf{6} \\ \hline
			\end{tabularx}
		\end{center}

	\section*{Evaluated criteria}
	\begin{enumerate}[label=\textbf{C\arabic*},start=5]
		\item Applies the Minimax Regret criterion by identifying the best payoff in each state, constructing the complete regret table, and selecting the alternative with the smallest maximum regret.
		\item Applies the Expected Monetary Value (EMV) or Bayes criterion by checking the state probabilities, calculating every alternative's expected payoff, and justifying the alternative with the highest EMV.
		\item Uses probability sensitivity analysis to determine remaining probabilities and feasible domains, derive EMV and threshold equations, and identify complete decision intervals including ties.
	\end{enumerate}

\sectionScored
	{Artisan Apiary Expansion}
	{C5-\sml{2}}
	{sec:garden-c5}
	{
	\label{sec:c5-exam-solutions}
	An artisan apiary must choose one of three hive-expansion plans. Annual net benefits are in thousands of dollars.
		\begin{center}
			\setlength{\tabcolsep}{3pt}
			\renewcommand{\arraystretch}{1.1}
			\begin{tabular}{lccc}
				\toprule
				& \multicolumn{3}{c}{Alternatives} \\
				& Large & Medium & Small \\ \midrule
				Excellent & 76 & 58 & 38 \\
				Average   & 34 & 49 & 43 \\
				Poor      & -14 & 23 & 35 \\ \bottomrule
			\end{tabular}
		\end{center}

	\QuestionsTasks{
		{5}/{subsec:garden-c5-solution}/
		{What decision would be selected using the Minimax Regret criterion? Justify your answer. \score{2}}
	}

	\phantomsection
	\label{subsec:garden-c5-solution}
	\subsection*{C5 - Solution \score{2}}

	Here $A_i$ follows the order of alternatives in the exam: Large, Medium, Small. States $S_j$ follow the exam order. All payoffs, regrets, and EMVs below are in thousands of dollars unless converted explicitly.

		The best payoff in each state is
		\[
		b(S_1)=\max\{76,58,38\}=76,\quad
		b(S_2)=\max\{34,49,43\}=49,\quad
		b(S_3)=\max\{-14,23,35\}=35.
		\]
		Using $r(A_i,S_j)=b(S_j)-P_{ay}(A_i,S_j)$ gives
		\begin{center}\renewcommand{\arraystretch}{1.25}\begin{tabular}{lcccc}\toprule
				& $S_1$ & $S_2$ & $S_3$ & Maximum regret\\ \midrule
				$A_1$&$76-76=0$&$49-34=15$&$35-(-14)=49$&$49$\\
				$A_2$&$76-58=18$&$49-49=0$&$35-23=12$&$18$\\
				$A_3$&$76-38=38$&$49-43=6$&$35-35=0$&$38$\\ \bottomrule\end{tabular}\end{center}
		Thus $\min\{49,18,38\}=18$, so the Minimax Regret criterion selects $\boxed{A_2}$, the medium expansion.

	\secInicio[sec:garden-c5][subsec:garden-c5-solution]
	}

\sectionScored
	{Regional Theater Capacity}
	{C5-\sml{2}}
	{sec:ferry-c5}
	{
	A theater operator is choosing a seating-capacity plan. Seasonal profit is in thousands of dollars; no reliable probabilities are available.
		\begin{center}\setlength{\tabcolsep}{2pt}\renewcommand{\arraystretch}{1.1}
			\begin{tabular}{lccc}\toprule
			Alternatives & High & Normal & Low\\ \midrule
			Add two seating sections&108&40&-28\\
			Add one seating section&82&61&14\\
			Lease a nearby hall&59&52&34\\
			Current seating&34&34&34\\ \bottomrule
			\end{tabular}
		\end{center}

	\QuestionsTasks{
		{5}/{subsec:ferry-c5-solution}/
		{What decision would be selected using the Minimax Regret criterion? Justify your answer. \score{2}}
	}

	\phantomsection
	\label{subsec:ferry-c5-solution}
	\subsection*{C5 - Solution \score{2}}

	Here $A_i$ follows the order of alternatives in the exam: add two seating sections, add one seating section, lease a nearby hall, current seating. States $S_j$ follow the exam order. All payoffs, regrets, and EMVs below are in thousands of dollars unless converted explicitly.

		The statewise best payoffs are $b(S_1)=108$, $b(S_2)=61$, and $b(S_3)=34$.
		\begin{center}\renewcommand{\arraystretch}{1.25}\begin{tabular}{lcccc}\toprule
				& $S_1$ & $S_2$ & $S_3$ & Maximum regret\\ \midrule
				$A_1$&$108-108=0$&$61-40=21$&$34-(-28)=62$&$62$\\
				$A_2$&$108-82=26$&$61-61=0$&$34-14=20$&$26$\\
				$A_3$&$108-59=49$&$61-52=9$&$34-34=0$&$49$\\
				$A_4$&$108-34=74$&$61-34=27$&$34-34=0$&$74$\\ \bottomrule\end{tabular}\end{center}
		Since $\min\{62,26,49,74\}=26$, the recommended decision is $\boxed{A_2}$: add one seating section.

	\secInicio[sec:ferry-c5][subsec:ferry-c5-solution]
	}

\sectionScored
	{Emergency Shelter Staffing}
	{C5-\sml{2}}
	{sec:clinic-c5}
	{
	A relief organization must select an emergency-shelter staffing plan. Net community value is in thousands of dollars.
		\begin{center}
			\setlength{\tabcolsep}{2pt}
			\renewcommand{\arraystretch}{1.1}
			\begin{tabular}{lccc}
				\toprule
				& \multicolumn{3}{c}{Alternatives} \\
				\cmidrule(lr){2-4}
				& Full & Flexible & Limited \\ \midrule
				Very high & 96 & 74 & 48 \\
				High      & 54 & 64 & 51 \\
				Moderate  & 24 & 49 & 46 \\
				Low       & -24 & 12 & 38 \\ \bottomrule
			\end{tabular}
		\end{center}

	\QuestionsTasks{
		{5}/{subsec:clinic-c5-solution}/
		{What decision would be selected using the Minimax Regret criterion? Justify your answer. \score{2}}
	}

	\phantomsection
	\label{subsec:clinic-c5-solution}
	\subsection*{C5 - Solution \score{2}}

	Here $A_i$ follows the order of alternatives in the exam: Full, Flexible, Limited. States $S_j$ follow the exam order. All payoffs, regrets, and EMVs below are in thousands of dollars unless converted explicitly.

		The best-payoff vector is
		\[
		(b(S_1),b(S_2),b(S_3),b(S_4))=(96,64,49,38).
		\]
		\begin{center}\renewcommand{\arraystretch}{1.25}\begin{tabular}{lccccc}\toprule
				&$S_1$&$S_2$&$S_3$&$S_4$&Maximum regret\\ \midrule
				$A_1$&0&10&25&62&62\\
				$A_2$&22&0&0&26&26\\
				$A_3$&48&13&3&0&48\\ \bottomrule\end{tabular}\end{center}
		For example, $r(A_1,S_4)=38-(-24)=62$ and $r(A_3,S_2)=64-51=13$. The smallest maximum regret is $\min\{62,26,48\}=26$, so select $\boxed{A_2}$, the flexible staffing plan.

	\secInicio[sec:clinic-c5][subsec:clinic-c5-solution]
	}

\sectionScored
	{Coffee Roastery Product Mix}
	{C6-\sml{2}}
	{sec:farm-c6}
	{
	\label{sec:c6-exam-solutions}
	A coffee roastery is selecting a product mix. Profits are in thousands of dollars.
		\begin{center}\setlength{\tabcolsep}{2pt}\renewcommand{\arraystretch}{1.1}
			\begin{tabular}{lccc}\toprule
			&Strong&Stable&Weak\\
			$A_i \qquad \quad P(S_j)$&0.45&0.35&0.20\\ \midrule
			Single-origin&86&44&-12\\
			Balanced&64&58&22\\
			Classic&42&47&37\\ \bottomrule
			\end{tabular}
		\end{center}

	\QuestionsTasks{
		{6}/{subsec:farm-c6-solution}/
		{What decision would be selected using the EMV or Bayes criterion? Justify your answer. \score{2}}
	}

	\phantomsection
	\label{subsec:farm-c6-solution}
	\subsection*{C6 - Solution \score{2}}

	Here $A_i$ follows the order of alternatives in the exam: Single-origin, Balanced, Classic. States $S_j$ follow the exam order. All payoffs, regrets, and EMVs below are in thousands of dollars unless converted explicitly.

		First, $0.45+0.35+0.20=1$. Then
		\begin{align*}
EMV(A_1)&=(0.45)(86)+(0.35)(44)+(0.20)(-12)=38.7+15.4-2.4=51.7,\\
			EMV(A_2)&=(0.45)(64)+(0.35)(58)+(0.20)(22)=28.8+20.3+4.4=53.5,\\
			EMV(A_3)&=(0.45)(42)+(0.35)(47)+(0.20)(37)=18.9+16.45+7.4=42.75.
		\end{align*}
		The highest EMV is $\max\{51.7,53.5,42.75\}=53.5$. Therefore, choose $\boxed{A_2}$, the balanced product mix, with expected profit $\$53{,}500$.

	\secInicio[sec:farm-c6][subsec:farm-c6-solution]
	}

\sectionScored
	{Aquarium Program Scale}
	{C6-\sml{2}}
	{sec:museum-c6}
	{
	An aquarium is deciding the scale of a temporary education program. Payoffs are in thousands of dollars.
		\begin{center}\setlength{\tabcolsep}{2pt}\renewcommand{\arraystretch}{1.1}
			\begin{tabular}{lccc}\toprule
			&High&Expected&Low\\
			$A_i \qquad \qquad P(S_j)$&0.30&0.40&0.30\\ \midrule
			Major&98&38&-24\\
			Standard&73&53&18\\
			Compact&52&47&34\\
			Regular tours only&27&27&27\\ \bottomrule
			\end{tabular}
		\end{center}

	\QuestionsTasks{
		{6}/{subsec:museum-c6-solution}/
		{What decision would be selected using the EMV or Bayes criterion? Justify your answer. \score{2}}
	}

	\phantomsection
	\label{subsec:museum-c6-solution}
	\subsection*{C6 - Solution \score{2}}

	Here $A_i$ follows the order of alternatives in the exam: Major, Standard, Compact, Regular tours only. States $S_j$ follow the exam order. All payoffs, regrets, and EMVs below are in thousands of dollars unless converted explicitly.

		The probabilities are valid because $0.30+0.40+0.30=1$.
		\begin{align*}
EMV(A_1)&=(0.30)(98)+(0.40)(38)+(0.30)(-24)=29.4+15.2-7.2=37.40,\\
			EMV(A_2)&=(0.30)(73)+(0.40)(53)+(0.30)(18)=21.9+21.2+5.4=48.50,\\
			EMV(A_3)&=(0.30)(52)+(0.40)(47)+(0.30)(34)=15.6+18.8+10.2=44.60,\\
			EMV(A_4)&=(0.30)(27)+(0.40)(27)+(0.30)(27)=27.00.
		\end{align*}
		Because $48.50$ is the highest EMV, the Bayes decision is $\boxed{A_2}$, the standard program, with expected payoff $\$48{,}500$.

	\secInicio[sec:museum-c6][subsec:museum-c6-solution]
	}

\sectionScored
	{Wind-Farm Maintenance Staffing}
	{C6-\sml{2}}
	{sec:solar-c6}
	{
	A wind-farm operator is choosing a maintenance staffing plan. Payoffs are in thousands of dollars.
		\begin{center}\setlength{\tabcolsep}{1.5pt}\renewcommand{\arraystretch}{1.1}
			\begin{tabular}{lcccc}\toprule
			&Surge&Strong&Steady&Slow\\
			$A_i \qquad \qquad P(S_j)$&0.20&0.30&0.35&0.15\\ \midrule
			Rapid hiring&105&64&24&-32\\
			Moderate hiring&79&68&48&12\\
			Current staff&53&55&51&42\\ \bottomrule
			\end{tabular}
		\end{center}

	\QuestionsTasks{
		{6}/{subsec:solar-c6-solution}/
		{What decision would be selected using the EMV or Bayes criterion? Justify your answer. \score{2}}
	}

	\phantomsection
	\label{subsec:solar-c6-solution}
	\subsection*{C6 - Solution \score{2}}

	Here $A_i$ follows the order of alternatives in the exam: Rapid hiring, Moderate hiring, Current staff. States $S_j$ follow the exam order. All payoffs, regrets, and EMVs below are in thousands of dollars unless converted explicitly.

		The probability check is $0.20+0.30+0.35+0.15=1$.
		\begin{align*}
			EMV(A_1)&=(0.20)(105)+(0.30)(64)+(0.35)(24)+(0.15)(-32)=21+19.2+8.4-4.8=43.8,\\
EMV(A_2)&=(0.20)(79)+(0.30)(68)+(0.35)(48)+(0.15)(12)=15.8+20.4+16.8+1.8=54.8,\\
			EMV(A_3)&=(0.20)(53)+(0.30)(55)+(0.35)(51)+(0.15)(42)=10.6+16.5+17.85+6.3=51.25.
		\end{align*}
		Thus $\max\{43.8,54.8,51.25\}=54.8$, so choose $\boxed{A_2}$, moderate hiring, with expected payoff $\$54{,}800$.

	\secInicio[sec:solar-c6][subsec:solar-c6-solution]
	}

\sectionScored
	{Concert Merchandise Order}
	{C7-\sml{2}}
	{sec:catering-c7}
	{
	\label{sec:c7-exam-solutions}
	A concert vendor must choose between a large and a small merchandise order. Payoffs are in thousands of dollars.
		\begin{center}\setlength{\tabcolsep}{4pt}\renewcommand{\arraystretch}{1.15}
			\begin{tabular}{lcc}\toprule
			$A_i$&High ($S_1$)&Low ($S_2$)\\ \midrule
			Large &80&20\\
			Small &50&40\\ \bottomrule
			\end{tabular}
		\end{center}

	\QuestionsTasks{
		{7}/{subsec:catering-c7-solution}/
		{Let $P(S_1)=p$. State the preferred alternative for every probability value ($0\le p\le1$). \score{2}}
	}

	\phantomsection
	\label{subsec:catering-c7-solution}
	\subsection*{C7 - Solution \score{2}}

	Here $A_i$ follows the order of alternatives in the exam: Large, Small. States $S_j$ follow the exam order. All payoffs, regrets, and EMVs below are in thousands of dollars unless converted explicitly.
		The listed states are assumed mutually exclusive and exhaustive, so their probabilities sum to one. Thus $P(S_2)=1-p$. Nonnegative probabilities require $p\ge0$ and $1-p\ge0$, giving the feasible domain $0\le p\le1$.

		\begin{align*}
EMV(A_1)&=80p+20(1-p)=20+60p,\\
			EMV(A_2)&=50p+40(1-p)=40+10p.
		\end{align*}
		At the threshold, $20+60p=40+10p$, so $50p=20$ and $p=\frac{2}{5}=0.4$. Since $A_2$ has the higher EMV at $p=0$ and $A_1$ has the higher EMV at $p=1$,
		\[
		\boxed{\begin{cases}
				A_2,&0\le p<\frac{2}{5},\\
				A_1\text{ and }A_2\text{ are tied},&p=\frac{2}{5},\\
				A_1,&\frac{2}{5}<p\le1.
		\end{cases}}
		\]

	\secInicio[sec:catering-c7][subsec:catering-c7-solution]
	}

\sectionScored
	{Meal-Kit Launch}
	{C7-\sml{2}}
	{sec:bike-c7}
	{
	A meal-kit company is comparing a bold launch with a measured launch. Payoffs are in thousands of dollars.
		\begin{center}\setlength{\tabcolsep}{2pt}\renewcommand{\arraystretch}{1.1}
			\begin{tabular}{lccc}\toprule
				&High&Moderate&Low\\ 				
				$A_i$&$S_1$&$S_2$&$S_3$\\ 
				\midrule
				Bold&60&30&0\\
				Measured&35&45&30\\ \bottomrule
			\end{tabular}
		\end{center}

	\QuestionsTasks{
		{7}/{subsec:bike-c7-solution}/
		{Suppose $P(S_1)=p$, $P(S_2)=\frac{p}{3}$. State the preferred alternative for every feasible probability value of $p$. \score{2}}
	}

	\phantomsection
	\label{subsec:bike-c7-solution}
	\subsection*{C7 - Solution \score{2}}

	Here $A_i$ follows the order of alternatives in the exam: Bold, Measured. States $S_j$ follow the exam order. All payoffs, regrets, and EMVs below are in thousands of dollars unless converted explicitly.
		The listed states are assumed mutually exclusive and exhaustive, so their probabilities sum to one.

		Because the three state probabilities must sum to one,
		\[
		P(S_3)=1-p-\frac{p}{3}=1-\frac{4p}{3}.
		\]
		Nonnegative probabilities require $p\ge0$ and $1-\frac{4p}{3}\ge0$, so the feasible range is $0\le p\le\frac{3}{4}$.
		\begin{align*}
EMV(A_1)&=60p+30\left(\frac{p}{3}\right)+0\left(1-\frac{4p}{3}\right)=70p,\\
			EMV(A_2)&=35p+45\left(\frac{p}{3}\right)+30\left(1-\frac{4p}{3}\right)=30+10p.
		\end{align*}
		Equating the EMVs gives $70p=30+10p$, hence
		\[
		60p=30\quad\Rightarrow\quad p=\frac{1}{2}.
		\]
		Therefore,
		\[
		\boxed{\begin{cases}
				A_2,&0\le p<\frac{1}{2},\\
				A_1\text{ and }A_2\text{ are tied},&p=\frac{1}{2},\\
				A_1,&\frac{1}{2}<p\le\frac{3}{4}.
		\end{cases}}
		\]

	\secInicio[sec:bike-c7][subsec:bike-c7-solution]
	}

\sectionScored
	{Publishing-Platform Bid}
	{C7-\sml{2}}
	{sec:rainwater-c7}
	{
	A publishing firm must choose between a standard platform bid and an experimental platform bid. Payoffs are in thousands of dollars.
		\begin{center}\setlength{\tabcolsep}{2pt}\renewcommand{\arraystretch}{1.1}
			\begin{tabular}{lccc}\toprule
				&Immediate&Revision&Delayed\\ 
				$A_i$&$S_1$&$S_2$&$S_3$\\ 
				\midrule
				Standard&30&70&70\\
				Experimental&50&50&30\\ \bottomrule
			\end{tabular}
		\end{center}

	\QuestionsTasks{
		{7}/{subsec:rainwater-c7-solution}/
		{Suppose $P(S_1)=3p$, $P(S_2)=p$. State the preferred alternative for every feasible probability value of $p$. \score{2}}
	}

	\phantomsection
	\label{subsec:rainwater-c7-solution}
	\subsection*{C7 - Solution \score{2}}

	Here $A_i$ follows the order of alternatives in the exam: Standard, Experimental. States $S_j$ follow the exam order. All payoffs, regrets, and EMVs below are in thousands of dollars unless converted explicitly.
		The listed states are assumed mutually exclusive and exhaustive, so their probabilities sum to one.

		The remaining probability is
		\[
		P(S_3)=1-3p-p=1-4p.
		\]
		The probability conditions $p\ge0$ and $1-4p\ge0$ give the feasible range $0\le p\le\frac{1}{4}$.
		\begin{align*}
EMV(A_1)&=30(3p)+70p+70(1-4p)=70-120p,\\
			EMV(A_2)&=50(3p)+50p+30(1-4p)=30+80p.
		\end{align*}
		At indifference,
		\[
		70-120p=30+80p
		\quad\Rightarrow\quad 40=200p
		\quad\Rightarrow\quad p=\frac{1}{5}=0.2.
		\]
		Here $A_1$ is preferred below the threshold, while the faster-growing EMV of $A_2$ makes it preferred above the threshold:
		\[
		\boxed{\begin{cases}
				A_1,&0\le p<\frac{1}{5},\\
				A_1\text{ and }A_2\text{ are tied},&p=\frac{1}{5},\\
				A_2,&\frac{1}{5}<p\le\frac{1}{4}.
		\end{cases}}
		\]

	\secInicio[sec:rainwater-c7][subsec:rainwater-c7-solution]
	}

\end{document}
